\section{Related work}
\label{sec:related}

\paragraph{Decoding and the tail.} Nucleus sampling \citep{holtzman2020curious} truncates
the unreliable tail of the next-token distribution; locally typical sampling
\citep{meister2023typical} targets human-like surprisal; min-$p$
\citep{nguyen2025minp} scales the truncation with the model's confidence and
is exactly the coherence floor we use; contrastive decoding
\citep{li2023contrastive} and plug-and-play steering \citep{dathathri2020pplm}
shape the distribution toward or away from a reference. All of these regulate
the tail. Our anti-probable sampler \emph{courts} it, inside an entropy band
and above a floor, and we find that this buys surface novelty only.

\paragraph{Repetition, self-reinforcement and entrainment.} The collapse of a
base model left to itself is a well-studied phenomenon. \citet{zhu-etal-2023-penalty}
name the \emph{self-reinforcement effect} (a repeated token becomes more
probable each time it repeats) and suppress it with a repetition penalty
restricted to a recent window; the graded, windowed penalty we call
\emph{habituation} is a decoding-time operator of the same family, and we
claim no novelty for it. \citet{guan-huang-2023-mitigating} attribute the
tendency to repeat to a learning bias of maximum-likelihood training and
correct it with self-contrastive training; \citet{xu-etal-2023-look} track the
distance between the current next-token distribution and past ones to detect
repetition and topic drift and steer decoding accordingly. At the mechanistic
level, \citet{niu-etal-2025-llama} show that language models assign higher
probability to tokens present in their context even when those tokens are
irrelevant (\emph{contextual entrainment}), and locate heads that carry it.
Our bare arm is entrainment on the model's own output; what this paper adds
is not a new remedy for repetition but a measurement of what the remedy
leaves untouched (a habituated stream is fluent and still not creative) and
of what a second, structural operator adds on top of it.

\paragraph{Evaluating creativity, and the judge.} \citet{nakajima-etal-2026-beyond}
show that the widely used Divergent Association Task ignores appropriateness
and propose to score novelty conditional on it; our three separate
dimensions (surprise, connection, coherence), reported without a composite,
follow the same logic of not letting novelty be bought with incoherence.
\citet{saakyan2026death} show with 8,618 expert annotations that
$\approx$91\% of top-quartile $n$-gram-novel expressions are not judged
creative; our factual-paraphrase escape mode and the sampler's idea-level
null are independent confirmations of the same gap. Infini-gram
\citep{liu2024infinigram} and Rusty-DAWG \citep{merrill2024rustydawg} make
verbatim novelty against a training corpus computable at scale; we use
infini-gram against the OLMo-2 corpus. On the instrument itself,
\citet{fein-etal-2026-litbench} find that the strongest off-the-shelf judge
they tested agrees with human preferences on creative-writing pairs only 73\%
of the time, and \citet{haldar-hockenmaier-2025-rating} document low
intra-rater reliability of LLM judges across runs. Both cautions apply to
this study. We answer the second by scoring every window with $k=5$
independent calls and using the median, and by reporting the resolution of
the instrument; we answer the first only partly, with a second judge family
and with human raters on a subset, and we treat the judge as the study's
main limitation (Section~\ref{sec:limitations}).

\paragraph{Loops and steering in story generation.} Outer loops that steer a
generator are common in story generation: SWAG \citep{pei-etal-2024-swag}
lets a second model choose the next action for the story model; Collective
Critics \citep{bae-kim-2024-collective} refine a plan and its expression with
several critic models; STORYTELLER \citep{li-etal-2025-storyteller} adds a
plot-planning structure to keep long stories coherent and cohesive. Against
this literature our contribution is not the idea of an outer loop but its
reduction: the operator we isolate is one injected sentence on a preserved
context, without a planner, a critic or a reward; the factorial ablation that
shows which part of the loop carries the effect; and the temporal
characterization (period, decay, timing) that a planning framework does not
provide.

\paragraph{Verified search.} FunSearch \citep{romeraparedes2024funsearch} and
AlphaEvolve \citep{novikov2025alphaevolve} pair a language model with a hard
evaluator and let selection do the work. Our bin-packing null (the
anti-probable operator does not separate from plain sampling under verified
evolution) and the loop results suggest that the structure of the generation
loop, not the sampling operator, is where to look next.

\paragraph{Reading the residual stream.} The logit lens is known to be
fragile in intermediate layers, and the tuned lens \citep{belrose2023tunedlens}
was proposed as a less biased alternative. We use the logit lens only for a
coarse commitment layer and otherwise work with mean-centered residual
geometry; the network section is descriptive, and its readings should be
repeated with a tuned lens before they are taken as mechanism.

\paragraph{Creative cognition.} Creative thought as coupling of default-mode
and executive networks \citep{beaty2016dynamics}, with salience regions
coupling first \citep{beaty2018robust}; creativity as connecting semantically
distant concepts \citep{kenett2018semantic}; incubation effects
\citep{sio2009incubation}; predictive processing \citep{clark2013whatever};
dreams as anti-overfitting noise \citep{hoel2021overfitted}; Boden's
novelty/surprise/value triad \citep{boden2004creative}; conceptual blending
\citep{fauconnier2002way}; novelty search \citep{stanley2015greatness}. The
reverie loop of Section~\ref{sec:loop} was an explicit engineering of the
first three into a text loop; its ablation is, in a sense, a test of which of
these ingredients a language model needs, and the answer (an interruption on
a clock) is humbler than the hypothesis.
